% chapters/02_why_models_can_learn_at_test_time.tex
% Version 2.0 — B11-B26 compliant rewrite
\chapter{为什么模型能在测试时学习？}
\label{ch:why_test_time}

\epigraph{``There is always a signal, even without the label. If a model can predict its own rotation, it has already aligned its feature space enough to guess the class.''}{--- \citeauthor{sun2020ttt}, 2020}

%─────────────────────────────────────────────
\section{本章想回答什么问题}
\label{sec:ttt_question}
%─────────────────────────────────────────────

上一章确立了``训练结束后仍可继续学习''的必要性。
但这立刻引发了一个核心技术问题：

\begin{quote}
\textit{在测试阶段，我们没有真实标签。
没有标签就无法计算监督损失，无法计算梯度。
那么模型到底依靠什么信号来调整自己的参数？}
\end{quote}

\citet{sun2020ttt} 的回答是：用\textbf{自监督学习（Self-Supervised Learning）}
从输入数据本身构造伪标签，由此产生梯度信号。
这就是 Test-Time Training（TTT）2020 的核心思想。

本章将：
\begin{enumerate}
  \item 给出 TTT 2020 的完整算法描述（\S\ref{sec:ttt_math}）
  \item 严格推导为什么自监督辅助头能间接改善主任务性能（\S\ref{sec:ttt_theory}）
  \item 讨论这种方法的工程代价和局限性（\S\ref{sec:ttt_limits}），
        为后续章节引入``架构内化更新''的动机做铺垫
\end{enumerate}

%─────────────────────────────────────────────
\section{最小例子：旋转预测作为自监督信号}
\label{sec:ttt_minimal}
%─────────────────────────────────────────────

\begin{toybox}[title=玩具问题 B：旋转预测实现域适应]
\textbf{训练阶段}：
一个图像分类器在清晰、正放的动物照片上训练。
训练时同时引入辅助任务：
把每张图分别旋转 $0^\circ, 90^\circ, 180^\circ, 270^\circ$，
让网络预测旋转角度（4 分类问题）。

\textbf{测试阶段}：
输入图像变成了模糊、倾斜的监控画面——
这是一种\textbf{分布偏移（Distribution Shift）}。
分类头的性能急剧下降。

\textbf{TTT 修正}：
对每张测试图像，先执行 $K$ 步旋转预测的辅助优化：
\begin{enumerate}
  \item 对测试图施加 4 种旋转，构造伪标签
        $\{(R_j(x_\text{test}), j)\}_{j=0}^{3}$
  \item 以旋转分类损失为目标，对特征提取器 $\theta_e$ 做 $K$ 步梯度更新
  \item 用更新后的提取器运行主分类头
\end{enumerate}

\textbf{关键洞见}：
正确预测旋转角度需要网络理解图像的``上下左右''——
这迫使特征提取器适应当前图像的分布特性，
从而间接改善主分类性能。
\end{toybox}

%─────────────────────────────────────────────
\section{直觉解释：双头共享表征架构}
\label{sec:why_test_time_intuition}
%─────────────────────────────────────────────

TTT 2020 的架构如下：

\begin{center}
\begin{tabular}{ll}
\toprule
\textbf{组件} & \textbf{角色} \\
\midrule
$f_{\theta_e}$ & 特征提取器（共享骨干，如 ResNet） \\
$f_{\theta_m}$ & 主任务头（分类/回归） \\
$f_{\theta_s}$ & 辅助自监督头（如旋转预测） \\
\bottomrule
\end{tabular}
\end{center}

训练时，三个组件联合训练，使得 $\theta_e$ 同时学习对主任务和辅助任务有用的表征。

测试时，\textbf{只更新 $\theta_e$}（通过辅助损失的梯度），
主任务头 $\theta_m$ 和辅助头 $\theta_s$ 的参数保持冻结。

这个机制之所以有效，是因为一个关键假设：
\begin{quote}
\textit{使表征适应辅助任务（旋转预测）所需的特征调整，
与使表征适应主任务（分类）所需的特征调整，
共享足够大的子空间。}
\end{quote}

如果这个假设成立，那么``为了更好地预测旋转''而修正 $\theta_e$，
会``顺带''改善分类性能。

%─────────────────────────────────────────────
\section{数学形式化：TTT 2020 算法}
\label{sec:ttt_math}
%─────────────────────────────────────────────

\subsection{预训练阶段}

在训练集 $\{(x_i, y_i)\}_{i=1}^N$ 上联合优化：
\begin{equation}
  \min_{\theta_e, \theta_m, \theta_s} \;
  \frac{1}{N} \sum_{i=1}^{N} \Big[
    \mathcal{L}_\text{main}\!\left(f_{\theta_m}(f_{\theta_e}(x_i)),\, y_i\right)
    + \alpha \cdot \mathcal{L}_\text{ss}\!\left(
      f_{\theta_s}(f_{\theta_e}(\hat{x}_i)),\, \hat{y}_i
    \right)
  \Big]
  \label{eq:ttt_pretrain}
\end{equation}
其中：
\begin{itemize}
  \item $\mathcal{L}_\text{main}$ 是主任务损失（如交叉熵）
  \item $\mathcal{L}_\text{ss}$ 是自监督辅助损失
  \item $\hat{x}_i = P(x_i)$ 是对输入施加的自监督变换
        （旋转、遮盖、打乱等）
  \item $\hat{y}_i$ 是变换的伪标签
        （在旋转预测中即旋转角度的类别索引）
  \item $\alpha > 0$ 是平衡主损失和辅助损失的超参数
\end{itemize}

\subsection{测试时适应阶段}

对每一个测试样本 $x_\text{test}$，执行以下步骤：

\textbf{Step 1：保存参数快照}
\begin{equation}
  \theta_e^{(0)} \leftarrow \theta_e^*
  \qquad (\text{从预训练得到的参数开始})
\end{equation}

\textbf{Step 2：构造自监督信号}

施加变换族 $\{P_j\}_{j=1}^{J}$，得到 $J$ 组伪标记数据：
\begin{equation}
  \{(P_j(x_\text{test}),\, j)\}_{j=1}^{J}
\end{equation}

\textbf{Step 3：$K$ 步梯度更新（仅更新 $\theta_e$）}
\begin{equation}
  \theta_e^{(k)} = \theta_e^{(k-1)}
  - \eta \nabla_{\theta_e} \mathcal{L}_\text{ss}\!\left(
    f_{\theta_s^*}\!(f_{\theta_e^{(k-1)}}(P_j(x_\text{test}))),\, j
  \right), \quad k = 1, \ldots, K
  \label{eq:ttt_update}
\end{equation}

\textbf{Step 4：用适应后的参数进行推断}
\begin{equation}
  \hat{y} = f_{\theta_m^*}\!(f_{\theta_e^{(K)}}(x_\text{test}))
  \label{eq:ttt_inference}
\end{equation}

\textbf{Step 5：参数重置}
\begin{equation}
  \theta_e \leftarrow \theta_e^* \qquad (\text{恢复为预训练参数})
  \label{eq:ttt_reset}
\end{equation}

\begin{remark}[参数重置的必要性]
步骤 5 至关重要。如果不重置 $\theta_e$，
对第 $n$ 个测试样本的适应会``污染''对第 $n+1$ 个样本的推断。
由于相邻测试样本可能来自不同的分布偏移模式，
累积的参数漂移会迅速导致性能崩溃。
\end{remark}

%─────────────────────────────────────────────
\section{为什么辅助任务能帮助主任务？}
\label{sec:ttt_theory}
%─────────────────────────────────────────────

直觉上，辅助任务与主任务之间存在``表征共享''关系。
我们可以用以下定理给出一个形式化的条件。

\begin{theorem}[辅助损失的间接改善条件]
\label{thm:ttt_indirect}
设 $\theta_e^*$ 是预训练后的特征提取器参数，
$\Delta = -\eta \nabla_{\theta_e} \mathcal{L}_\text{ss}$
是自监督损失的一步更新。

如果主任务损失的梯度与辅助损失的梯度之间满足
\textbf{正相关条件}：
\begin{equation}
  \langle
    \nabla_{\theta_e} \mathcal{L}_\text{main}(\theta_e^*),\;
    \nabla_{\theta_e} \mathcal{L}_\text{ss}(\theta_e^*)
  \rangle > 0
  \label{eq:gradient_alignment}
\end{equation}
则辅助更新 $\theta_e^* + \Delta$ 在一阶近似下也会减少主任务损失：
\begin{equation}
  \mathcal{L}_\text{main}(\theta_e^* + \Delta)
  \approx \mathcal{L}_\text{main}(\theta_e^*)
  - \eta \langle
    \nabla \mathcal{L}_\text{main},\, \nabla \mathcal{L}_\text{ss}
  \rangle < \mathcal{L}_\text{main}(\theta_e^*)
\end{equation}
\end{theorem}

\begin{proof}
一阶 Taylor 展开：
$\mathcal{L}_\text{main}(\theta_e^* + \Delta)
\approx \mathcal{L}_\text{main}(\theta_e^*)
+ \langle \nabla \mathcal{L}_\text{main}, \Delta \rangle$。
代入 $\Delta = -\eta \nabla \mathcal{L}_\text{ss}$
即得 $\approx \mathcal{L}_\text{main}(\theta_e^*)
- \eta \langle \nabla \mathcal{L}_\text{main},
\nabla \mathcal{L}_\text{ss} \rangle$。
由条件~\eqref{eq:gradient_alignment}，
括号内的内积 $> 0$，因此主损失下降。
\end{proof}

\begin{remark}
条件~\eqref{eq:gradient_alignment} 并不总是成立。
当辅助任务与主任务之间的特征共享不足时
（例如用``颜色直方图匹配''辅助任务来帮助``语义分割''），
梯度可能正交甚至反向，
导致测试时适应反而损害主任务性能。
这是 TTT 方法的核心风险。
\end{remark}

%─────────────────────────────────────────────
\section{计算代价分析与工程局限}
\label{sec:ttt_limits}
%─────────────────────────────────────────────

TTT 2020 的主要工程问题在于测试时的\textbf{计算开销}。

\begin{center}
\begin{tabular}{p{3.5cm} p{4.5cm} p{4.5cm}}
\toprule
& \textbf{标准推断} & \textbf{TTT 推断} \\
\midrule
每样本计算量 & 1 次前向传播 & $K \times J$ 次前向 + $K$ 次反向 \\
GPU 内存 & 仅存激活值 & 需存梯度计算图 \\
延迟 & $\sim$1ms & $\sim$50--500ms (取决于 $K$) \\
参数管理 & 无 & 需保存/恢复 $\theta_e^*$ \\
\bottomrule
\end{tabular}
\end{center}

对于\textbf{逐 token 的序列模型}（如语言模型），
TTT 2020 的``全局反向传播 + 参数重置''范式完全不可行：
每生成一个 token 就要对整个网络反向传播一次。

\begin{insight}[title=从外部微调到架构内化]
TTT 2020 的核心贡献是证明了
``测试时学习''的\textbf{可行性}——模型确实可以利用自监督信号适应。
但它的工程瓶颈指向了后续研究的方向：
\textbf{能否把``测试时的参数更新''
内化到网络架构本身中}——
不是在外部跑反向传播，
而是设计一种层结构，
其前向传播天然等价于执行参数更新？
这正是第~\ref{ch:bilevel}--\ref{ch:ttt_layers}~章的主题。
\end{insight}

%─────────────────────────────────────────────
\section{神经科学直觉与类比边界}
\label{sec:ttt_neuro}
%─────────────────────────────────────────────

\begin{neurosciencebox}[title=神经科学直觉：感觉-运动的在线校准]

\textbf{类比对象}：
人类的\textbf{棱镜适应（Prism Adaptation）}现象
可以帮助理解 TTT 的核心机制。

当受试者佩戴会使视觉偏移 $20^\circ$ 的棱镜时，
初始的手眼协调完全失准（类似分布偏移后模型性能下降）。
但仅需数次伸手抓握尝试，受试者就能重新校准——
\textbf{即使没有人告诉他``正确的手臂位置在哪里''}。
校准依赖的是\textbf{自身的本体感觉反馈}
（手臂的运动感知信号），而非外部标签。

\textbf{共同之处（算法层面）}：
\begin{center}
\begin{tabular}{p{5cm} p{5cm}}
\toprule
\textbf{棱镜适应} & \textbf{TTT 2020} \\
\midrule
视觉偏移 = 分布偏移 & 测试图像来自新分布 \\
本体感觉反馈 = 自监督信号 & 旋转预测提供伪标签 \\
小脑校准视觉-运动映射 & 梯度更新特征提取器 \\
摘除棱镜后需重新适应 & 参数重置避免跨样本污染 \\
\bottomrule
\end{tabular}
\end{center}

\end{neurosciencebox}

\begin{analogyboundarybox}[title=这个类比在哪里失效]

\textbf{一、更新规则不同（实现机制层）}

\begin{tabular}{p{5cm} p{8cm}}
\textbf{TTT 2020}
& 通过显式的全局反向传播计算梯度，
  沿着整个计算图将误差传播到 $\theta_e$ 的每个参数。
  需要存储完整的计算图。 \\
\textbf{棱镜适应}
& 小脑的误差校正依赖于\textbf{局部的}
  攀缘纤维信号和 Purkinje 细胞的突触可塑性，
  是一种逐层局部的适应，不存在全局误差反传。 \\
\end{tabular}

\medskip\textbf{二、适应范围不同（表征层）}

TTT 更新的是神经网络中\textbf{所有共享层}的参数——
这是一种全局性的表征调整。
生物的棱镜适应主要局限于
\textbf{小脑-顶叶}回路的特定连接，
不会改写整个视觉皮层的表征。

\medskip\textbf{读者不应从这个类比推出}：
\begin{itemize}
  \item \textbf{[X]} ``TTT 2020 实现了类脑的在线适应''
  \item \textbf{[X]} ``反向传播等价于小脑的误差校正''
  \item \textbf{[X]} ``参数重置等价于摘除棱镜后的重新适应''
\end{itemize}

\medskip
\textbf{类比成立的层次}：
\textit{算法层面}——两者都使用非任务本身的辅助信号
来间接校准任务相关的表征。

\textbf{类比失效的层次}：
实现机制（全局反向传播 vs.\ 局部突触可塑性）、
更新范围（所有共享层 vs.\ 特定回路）均不同。

\end{analogyboundarybox}

%─────────────────────────────────────────────
\section{代表论文精读}
\label{sec:why_test_time_papers}
%─────────────────────────────────────────────

\paragraph{\citet{sun2020ttt} — Test-Time Training with Self-Supervision for Generalization under Distribution Shifts}
这是 TTT 范式的奠基论文。
核心贡献：证明了旋转预测作为自监督辅助任务，
能够有效改善分布偏移下的分类性能。
在 CIFAR-10-C、ImageNet-C 等基准上，
TTT 显著优于不做任何适应的基线。
更重要的是，该论文建立了``测试时利用自监督信号''
这一范式的可行性证明（existence proof）——
为后续所有测试时学习研究奠定了基础。

\textit{本教材用处}：证明无标签条件下的测试时适应是可行的。
本章的整个框架（双头架构、辅助损失、参数重置）
直接来源于这篇论文。

\paragraph{\citet{krause2018dynamiceval} — Dynamic Evaluation of Neural Sequence Models}
在语言模型领域，Dynamic Evaluation 是一种更早期的测试时适应策略：
在部署时持续用语言模型的自回归损失（下一个 token 预测）
对模型参数进行在线微调。
这可以被理解为一种不需要额外辅助头的 TTT 方案，
因为语言模型的主任务本身就是自监督的。

\textit{本教材用处}：说明 TTT 思想在 NLP 领域的自然对应物，
并为后续的 TTT-Layer（第~\ref{ch:ttt_layers}~章）提供额外动机。

%─────────────────────────────────────────────
\section{和前后章节的关系}
\label{sec:why_test_time_bridges}
%─────────────────────────────────────────────

\begin{itemize}
  \item \textbf{来自第~\ref{ch:intro}~章}：
    上一章确立了``测试阶段学习''的必要性。
    本章给出了最简单的实现方案——
    外部挂载自监督辅助头 + 反向传播。

  \item \textbf{通向第~\ref{ch:bilevel}~章}：
    TTT 2020 的``预训练阶段学习辅助头''和``测试阶段用辅助头微调''
    天然形成了一个双层优化结构：
    外循环（预训练）学习``如何设定辅助任务的参数''，
    内循环（测试时）用辅助任务微调特征提取器。
    这正是 MAML 等元学习方法的核心框架。

  \item \textbf{通向第~\ref{ch:fastweights}--\ref{ch:ttt_layers}~章}：
    本章揭示的工程瓶颈——
    ``每次推断都需要全局反向传播''——
    驱动了后续研究的核心目标：
    把测试时更新\textbf{内化到网络层中}，
    使前向传播本身就执行参数适应，
    无需外部反向传播。
\end{itemize}

%─────────────────────────────────────────────
\section{实验}
\label{sec:ttt_lab}
%─────────────────────────────────────────────

\begin{labbox}[title=Lab 2：旋转自监督 TTT 的完整实现]
\textbf{目标}：在 CIFAR-10 上实现 TTT 2020，
并测量适应效果和计算开销。

\begin{enumerate}
  \item \textbf{构建双头网络}：
    使用 ResNet-18 作为特征提取器 $\theta_e$，
    接一个 10 类分类头 $\theta_m$ 和一个 4 类旋转预测头 $\theta_s$。

  \item \textbf{联合预训练}：
    在 CIFAR-10 训练集上用式~\eqref{eq:ttt_pretrain} 联合训练，
    $\alpha = 0.5$。记录收敛后的测试准确率（无分布偏移）。

  \item \textbf{制造分布偏移}：
    对 CIFAR-10 测试集施加高斯模糊 + 随机旋转（非 $90^\circ$ 倍数），
    记录不做 TTT 时的准确率下降。

  \item \textbf{实施 TTT}：
    对每张测试图执行式 \eqref{eq:ttt_update}，$K=1,3,10$。
    记录准确率恢复情况。

  \item \textbf{测量计算开销}：
    比较``标准推断''和``TTT 推断''的每张图延迟和 GPU 内存占用。

  \item \textbf{不重置实验（对照组）}：
    去掉式~\eqref{eq:ttt_reset} 的参数重置，
    观察在处理 100 张连续测试图后的性能衰退曲线。
\end{enumerate}

\textbf{预期结果}：
\begin{itemize}
  \item TTT ($K=3$) 应将准确率恢复 5--15\%
  \item 不重置参数会导致准确率在约 50 张图后开始下降
  \item TTT 每张图的延迟约为标准推断的 10--50 倍
\end{itemize}
\end{labbox}

%─────────────────────────────────────────────
\section{常见误解}
\label{sec:ttt_misconceptions}
%─────────────────────────────────────────────

\begin{misconceptionbox}
\textbf{误解 1}：只要设计任意一个自监督辅助任务，
TTT 就能改善主任务性能。

\textbf{纠正}：
只有当辅助任务的梯度与主任务的梯度之间
满足正相关条件~\eqref{eq:gradient_alignment} 时，
辅助更新才能帮助主任务。
如果辅助任务所需的特征调整与主任务正交或冲突，
TTT 不仅无效，还可能有害。
选择合适的辅助任务是 TTT 方法的核心设计难点。
\end{misconceptionbox}

\begin{misconceptionbox}
\textbf{误解 2}：既然我们已经有了 TTT 2020，
为什么还需要设计特殊的网络层（TTT-Layer 等）？
直接给任何模型挂一个辅助头不就行了？

\textbf{纠正}：
TTT 2020 对每个测试样本都需要执行\textbf{完整的反向传播}，
计算开销远超正常推断。对于序列模型（每个 token 都需要适应），
这种开销是灾难性的。后续章节的核心动机正是：
把``测试时更新''从外部的反向传播
\textbf{内化}为前向传播层的固有行为，
使推断效率与普通层相当。
这是架构层面的革命，不是算法层面的修补。
\end{misconceptionbox}

%─────────────────────────────────────────────
\section{练习题}
\label{sec:why_test_time_exercises}
%─────────────────────────────────────────────

\begin{exercise}
\textbf{参数污染分析}：
如果在式~\eqref{eq:ttt_update} 的 $K$ 步更新后
不执行参数重置（式~\eqref{eq:ttt_reset}），
而是直接用 $\theta_e^{(K)}$ 处理下一张测试图。
假设第二张测试图来自与第一张完全不同的分布偏移，
用一阶 Taylor 展开分析为什么 $\theta_e^{(K)}$
对第二张图的主任务损失可能比 $\theta_e^*$ 更差。
\end{exercise}

\begin{exercise}
\textbf{梯度对齐条件}：
设主任务是 CIFAR-10 物体分类，考虑以下三种辅助任务：
（a）旋转预测、（b）颜色通道排列预测、（c）图像拼图重组。
分析每种辅助任务在``含高斯模糊''和``含几何变形''
两种分布偏移下，梯度对齐条件~\eqref{eq:gradient_alignment}
是否可能被满足，并说明理由。
\end{exercise}

\begin{exercise}
\textbf{从 CV 到 NLP}：
为基于 Transformer 的语言模型设计一种自监督辅助任务，
使其能在测试时适应领域偏移
（例如从学术论文切换到社交媒体文本）。
你的设计需要说明：
（a）辅助任务的具体定义、
（b）伪标签如何构造、
（c）最可能满足的梯度对齐条件。
\end{exercise}

\begin{exercise}
\textbf{步数 $K$ 的权衡}：
分析 TTT 更新步数 $K$ 的增加对以下指标的影响：
（a）测试时计算开销、
（b）辅助任务损失 $\mathcal{L}_\text{ss}$、
（c）主任务性能。
说明为什么存在一个最优的 $K^*$，
超过该步数后主任务性能可能反而下降（过拟合于单样本）。
\end{exercise}

%─────────────────────────────────────────────
\section{本章小结}
\label{sec:why_test_time_summary}
%─────────────────────────────────────────────

\begin{itemize}
  \item \textbf{TTT 2020 的核心思想}：
    通过自监督辅助任务（如旋转预测），
    在测试时为模型提供无需真实标签的梯度信号。
    辅助更新调整共享特征提取器，间接改善主任务性能。

  \item \textbf{有效性条件}：
    辅助损失与主损失的梯度之间必须满足正相关条件
    （定理~\ref{thm:ttt_indirect}），
    否则 TTT 可能无效甚至有害。

  \item \textbf{工程代价}：
    每个测试样本需要 $K$ 次完整的反向传播，
    计算开销远超标准推断。
    参数重置是必要的，否则会导致跨样本污染。

  \item \textbf{通向后续章节的桥梁}：
    TTT 2020 的工程瓶颈
    （全局反向传播、高延迟、参数管理复杂）
    直接驱动了``把测试时更新内化到网络架构中''的研究方向——
    这正是从第~\ref{ch:bilevel}~章开始的主线。

  \item \textbf{历史地位}：
    TTT 2020 是``测试时学习''的\textbf{可行性证明}——
    它证明了这条路走得通。
    后续的快权重、TTT-Layer、Titans 等是
    ``如何高效走这条路''的进化。
\end{itemize}

\paragraph{读完本章，你应该能够：}
\begin{itemize}
  \item 写出 TTT 2020 的完整算法（预训练 + 测试时更新 + 推断 + 重置）
  \item 用一阶 Taylor 展开论证辅助梯度何时帮助/何时损害主任务
  \item 解释为什么参数重置是必要的
  \item 量化 TTT 的计算开销，理解其对序列模型不可行的原因
  \item 阐述为什么 TTT 2020 的工程瓶颈
        引出了``内化更新到网络层''的后续研究方向
\end{itemize}
